\documentclass[10pt,a4paper]{amsart}
\usepackage{amsmath,amssymb,amsthm,graphicx}
\usepackage{verbatim}
\newenvironment{reference}{\comment}{\endcomment}
%\newenvironment{reference}{}{}
\newenvironment{slogan}{\comment}{\endcomment}
\newenvironment{history}{\comment}{\endcomment}

% For commutative diagrams we use Xy-pic
\usepackage[all]{xy}

% We use 2cell for 2-commutative diagrams.
\xyoption{2cell}
% Theorem environments.
%
\theoremstyle{plain}
\newtheorem{theorem}[subsection]{Theorem}
\newtheorem{proposition}[subsection]{Proposition}
\newtheorem{lemma}[subsection]{Lemma}

\theoremstyle{definition}
\newtheorem{definition}[subsection]{Definition}
\newtheorem{example}[subsection]{Example}
\newtheorem{exercise}[subsection]{Exercise}
\newtheorem{situation}[subsection]{Situation}

\theoremstyle{remark}
\newtheorem{remark}[subsection]{Remark}
\newtheorem{remarks}[subsection]{Remarks}

\numberwithin{equation}{subsection}

% Macros
%
\def\lim{\mathop{\mathrm{lim}}\nolimits}
\def\colim{\mathop{\mathrm{colim}}\nolimits}
\def\Spec{\mathop{\mathrm{Spec}}}
\def\Hom{\mathop{\mathrm{Hom}}\nolimits}
\def\Ext{\mathop{\mathrm{Ext}}\nolimits}
\def\SheafHom{\mathop{\mathcal{H}\!\mathit{om}}\nolimits}
\def\SheafExt{\mathop{\mathcal{E}\!\mathit{xt}}\nolimits}
\def\Sch{\mathit{Sch}}
\def\Mor{\mathop{\mathrm{Mor}}\nolimits}
\def\Ob{\mathop{\mathrm{Ob}}\nolimits}
\def\Sh{\mathop{\mathit{Sh}}\nolimits}
\def\NL{\mathop{N\!L}\nolimits}
\def\CH{\mathop{\mathrm{CH}}\nolimits}
\def\proetale{{pro\text{-}\acute{e}tale}}
\def\etale{{\acute{e}tale}}
\def\QCoh{\mathit{QCoh}}
\def\Ker{\mathop{\mathrm{Ker}}}
\def\Im{\mathop{\mathrm{Im}}}
\def\Coker{\mathop{\mathrm{Coker}}}
\def\Coim{\mathop{\mathrm{Coim}}}

% Boxtimes
%
\DeclareMathSymbol{\boxtimes}{\mathbin}{AMSa}{"02}

%
% Macros for moduli stacks/spaces
%
\def\QCohstack{\mathcal{QC}\!\mathit{oh}}
\def\Cohstack{\mathcal{C}\!\mathit{oh}}
\def\Spacesstack{\mathcal{S}\!\mathit{paces}}
\def\Quotfunctor{\mathrm{Quot}}
\def\Hilbfunctor{\mathrm{Hilb}}
\def\Curvesstack{\mathcal{C}\!\mathit{urves}}
\def\Polarizedstack{\mathcal{P}\!\mathit{olarized}}
\def\Complexesstack{\mathcal{C}\!\mathit{omplexes}}
% \Pic is the operator that assigns to X its picard group, usage \Pic(X)
% \Picardstack_{X/B} denotes the Picard stack of X over B
% \Picardfunctor_{X/B} denotes the Picard functor of X over B
\def\Pic{\mathop{\mathrm{Pic}}\nolimits}
\def\Picardstack{\mathcal{P}\!\mathit{ic}}
\def\Picardfunctor{\mathrm{Pic}}
\def\Deformationcategory{\mathcal{D}\!\mathit{ef}}
\usepackage{cleveref}
\crefname{originaldefinition}{original definition}{original definitions}
\Crefname{originaldefinition}{Original definition}{Original definitions}
\makeatletter
\newlabel{formal-spaces-definition-formal-algebraic-space}{{0AIM}{source}{}{sourcepointer-1}{}}
\newlabel{formal-spaces-definition-formal-algebraic-space@cref}{{[originaldefinition][0][]0AIM}{[1][0][]source}{}{}{}}
\makeatother
\crefname{originallemma}{original lemma}{original lemmas}
\Crefname{originallemma}{Original lemma}{Original lemmas}
\makeatletter
\newlabel{formal-spaces-lemma-formal-completion-types}{{0AQ1}{source}{}{sourcepointer-2}{}}
\newlabel{formal-spaces-lemma-formal-completion-types@cref}{{[originallemma][0][]0AQ1}{[1][0][]source}{}{}{}}
\makeatother
\crefname{originalreference}{original reference}{original references}
\Crefname{originalreference}{Original reference}{Original references}
\makeatletter
\newlabel{lemma-rig-smooth}{{0GAK}{source}{}{sourcepointer-3}{}}
\newlabel{lemma-rig-smooth@cref}{{[originalreference][0][]0GAK}{[1][0][]source}{}{}{}}
\makeatother
\crefname{originalreference}{original reference}{original references}
\Crefname{originalreference}{Original reference}{Original references}
\makeatletter
\newlabel{lemma-rig-smooth-morphisms}{{0GCQ}{source}{}{sourcepointer-4}{}}
\newlabel{lemma-rig-smooth-morphisms@cref}{{[originalreference][0][]0GCQ}{[1][0][]source}{}{}{}}
\makeatother
\crefname{originallemma}{original lemma}{original lemmas}
\Crefname{originallemma}{Original lemma}{Original lemmas}
\makeatletter
\newlabel{spaces-properties-lemma-cover-by-union-affines}{{03FX}{source}{}{sourcepointer-5}{}}
\newlabel{spaces-properties-lemma-cover-by-union-affines@cref}{{[originallemma][0][]03FX}{[1][0][]source}{}{}{}}
\makeatother
\title{Completion along a closed subset of a modification smooth over the complement is rig-smooth}
\author{The Stacks Project Authors}
\date{stacks-project:restricted.tex:local-snapshot-20261012; GNU FDL 1.2 or later}
\begin{document}
\maketitle
\noindent\textit{Editorial note.} All delivered mathematical text is the author's own, copied from the pinned Stacks Project chapter file restricted.tex; no statement or proof step was rewritten, reordered or removed.\par\medskip
\noindent\textit{Editorial note.} The counted proof is the author's complete argument: a double etale affine covering reduction, a locality argument for rig-smooth morphisms of formal algebraic spaces, and an affine algebra conclusion via the I-adic completion criterion, all inside the delivered span\par\medskip
\noindent\textit{Editorial note.} The label lemma-smooth-gives-rig-smooth names the counted statement, and the proof environment immediately following it in the source is delivered verbatim as its proof.\par\medskip
\noindent\textit{Editorial note.} The counted statement is carried by the author's own lemma environment in the chapter source.\par\medskip
\noindent\textit{Editorial note.} The proof opens with the author's sentence 'We urge the reader to skip the details.' and then nevertheless delivers the complete reduction and affine argument; per the quality criterion this reading suggestion is not an omission, and the verification is delivered verbatim.\par\medskip
\noindent\textit{Allowed standard prerequisites (editorial).}\begin{itemize}
\item The modification setting: f : X' to X with U' to U smooth over the complement of T, and cartesianity of the diagrams formed by completing along inverse images of T
\item The covering definition of locally Noetherian formal algebraic spaces, and rig-smoothness being local on the source and target for morphisms of formal algebraic spaces
\item The affine algebra criterion identifying rig-smoothness of completed finite type ring maps
\end{itemize}\medskip
\ifdefined\hypertarget\hypertarget{sourcepointer-1}{}\fi
\noindent\textit{Source reference.} Original definition 0AIM: sources/stacks-chow/tags.txt, source lines 13270--13270. Prerequisite pointer invoked at proof line(s) 6151; it names a standard fact of a sibling chapter that the counted argument uses, and every step of the counted proof built on it is retained verbatim in the delivered body.\par\medskip
\ifdefined\hypertarget\hypertarget{sourcepointer-2}{}\fi
\noindent\textit{Source reference.} Original lemma 0AQ1: sources/stacks-chow/tags.txt, source lines 13494--13494. Prerequisite pointer invoked at proof line(s) 6177; it names a standard fact of a sibling chapter that the counted argument uses, and every step of the counted proof built on it is retained verbatim in the delivered body.\par\medskip
\ifdefined\hypertarget\hypertarget{sourcepointer-3}{}\fi
\noindent\textit{Source reference.} Original reference 0GAK: sources/stacks-restricted/pinned-original-source.tex, source lines 943--943. Prerequisite pointer invoked at proof line(s) 6132; it names a standard fact of the same chapter that the counted argument uses, and every step of the counted proof built on it is retained verbatim in the delivered body.\par\medskip
\ifdefined\hypertarget\hypertarget{sourcepointer-4}{}\fi
\noindent\textit{Source reference.} Original reference 0GCQ: sources/stacks-restricted/pinned-original-source.tex, source lines 4642--4642. Prerequisite pointer invoked at proof line(s) 6152; it names a standard fact of the same chapter that the counted argument uses, and every step of the counted proof built on it is retained verbatim in the delivered body.\par\medskip
\ifdefined\hypertarget\hypertarget{sourcepointer-5}{}\fi
\noindent\textit{Source reference.} Original lemma 03FX: sources/stacks-chow/tags.txt, source lines 4508--4508. Prerequisite pointer invoked at proof line(s) 6137; it names a standard fact of a sibling chapter that the counted argument uses, and every step of the counted proof built on it is retained verbatim in the delivered body.\par\medskip
\begin{lemma}
\label{lemma-smooth-gives-rig-smooth}
In the situation above. If $X$ is locally Noetherian,
$f$ is locally of finite type, and $U' \to U$ is smooth, then
$f_{/T}$ is rig-smooth.
\end{lemma}

\begin{proof}
The strategy of the proof is this: reduce to the case where $X$ and $X'$
are affine, translate the affine case into algebra, and finally apply
Lemma \ref{lemma-rig-smooth}. We urge the reader to skip the details.

\medskip\noindent
Choose a surjective \'etale morphism $W \to X$ with $W = \coprod W_i$
a disjoint union of affine schemes, see Properties of Spaces, Lemma
\ref{spaces-properties-lemma-cover-by-union-affines}.
For each $i$ choose a surjective \'etale morphism $W'_i \to W_i \times_X X'$
where $W'_i = \coprod W'_{ij}$ is a disjoint union of affines.
In particular $\coprod W'_{ij} \to X'$ is surjective and \'etale.
Denote $f_{ij} : W_{ij} \to W_i$ the given morphism.
Denote $T_i \subset W_i$ and $T'_{ij} \subset W_{ij}$ the inverse
images of $T$. Since taking the completion along the inverse
image of $T$ produces cartesian diagrams (see above) we
have $(W_i)_{/T_i} = W_i \times_X X_{/T}$ and similarly
$(W'_{ij})_{/T'_{ij}} = W'_{ij} \times_{X'} X'_{/T'}$.
Moreover, recall that $(W_i)_{/T_i}$ and $(W'_{ij})_{/T'_{ij}}$
are affine formal algebraic spaces.
Hence $\{W'_{ij})_{/T'_{ij}} \to X'_{/T'}\}$ is a covering
as in Formal Spaces, Definition
\ref{formal-spaces-definition-formal-algebraic-space}.
By Lemma \ref{lemma-rig-smooth-morphisms}
we see that it suffices to prove that
$$
(W'_{ij})_{/T'_{ij}} \longrightarrow (W_i)_{/T_i}
$$
is rig-smooth. Observe that $W'_{ij} \to W_i$
is locally of finite type and induces a smooth morphism
$W'_{ij} \setminus T'_{ij} \to W_i \setminus T_i$
(as this is true for $f$ and these properties of morphisms are
\'etale local on the source and target).
Observe that $W_i$ is locally Noetherian (as $X$ is locally Noetherian
and this property is \'etale local on the algebraic space).
Hence it suffices to prove the lemma when
$X$ and $X'$ are affine schemes.

\medskip\noindent
Assume $X = \Spec(A)$ and $X' = \Spec(A')$ are affine schemes.
Since $X$ is Noetherian, we see that $A$ is Noetherian.
The morphism $f$ is given by a ring map $A \to A'$ of finite type.
Let $I \subset A$ be an ideal cutting out $T$. Then $IA'$ cuts out $T'$.
Also $\Spec(A') \to \Spec(A)$ is smooth over $\Spec(A) \setminus T$.
Let $A^\wedge$ and $(A')^\wedge$ be the $I$-adic
completions. We have $X_{/T} = \text{Spf}(A^\wedge)$ and
$X'_{/T'} = \text{Spf}((A')^\wedge)$, see proof of
Formal Spaces, Lemma
\ref{formal-spaces-lemma-formal-completion-types}.
By Lemma \ref{lemma-rig-smooth} we see that
$(A')^\wedge$ is rig-smooth over $(A. I)$
which in turn means that $A^\wedge \to (A')^\wedge$ is rig-smooth which
finally implies that $X'_{/T'} \to X_{/T}$ is rig smooth by
Lemma \ref{lemma-rig-smooth-morphisms}.
\end{proof}

\end{document}
